\subsection{Cyclic correlation and preservation of path order}
\label{sec:rec27-memoria}

The record \eqref{act21:tpk:archivo} preserves the ordered
concatenation of samples. The following example isolates order
information extracted by two distinct readers from a twelve-position
word, retaining the domain of each reading.

Let $\tau=(\tau_1,\ldots,\tau_{12})\in\{-1,0,1\}^{12}$ be an ordered
word, with indices taken modulo $12$. Define
\[
\begin{aligned}
 G_a&=\{a,a+4,a+8\}\quad(1\leq a\leq4),\\
 \nu_{120}(\tau)&=\sum_{a=1}^4\operatorname{sgn}
 \left(\sum_{m\in G_a}\tau_m\right),\\
 \nu_{270}(\tau)&=\sum_{m=1}^{12}\tau_m\tau_{m+3}.
\end{aligned}
\]
The marginal histogram is $p_j(\tau)=\#\{m:\tau_m=j\}/12$ for
$j\in\{-1,0,1\}$, and its entropy is
$H(\tau)=-\sum_j p_j(\tau)\log p_j(\tau)$, with $0\log0=0$.

\begin{proposition}[Order information beyond the marginal distribution]
There exist words with identical histograms and distinct correlations
$\nu_{270}$, even when $\nu_{120}$ has the same value.
\end{proposition}
\begin{proof}
Consider
\[
 \tau^{(a)}=(1,1,0,0,0,0,0,0,0,0,0,0),\qquad
 \tau^{(b)}=(1,0,0,1,0,0,0,0,0,0,0,0).
\]
Both have $p_1=1/6$, $p_0=5/6$ and $p_{-1}=0$, so their marginal
entropies coincide. In each case, two of the sets $G_a$ contain a unit and
the other two contain only zeros: $\nu_{120}=2$.
In $\tau^{(a)}$, no pair of units has cyclic separation three, whereas in
$\tau^{(b)}$ precisely the term $m=1$ contributes.
Thus $\nu_{270}(\tau^{(a)})=0$ and $\nu_{270}(\tau^{(b)})=1$.
\end{proof}

This separation identifies the information retained by the ordered record:
incidences between positions in addition to symbol frequencies.
A reduced state retaining only the histogram assigns the same output to
both words and loses this correlation. Marginal entropy, uncertainty
between bases and entanglement entropy are defined on different objects;
relating them requires retaining the distribution, the change of basis,
and the state together with its partition, respectively.
